\section{Original theta Hankel witness: independent proof and full packet transport}\label{C47stage:WBR:original-theta-hankel-witness-independent-proof-and-full-packet-transport}

This is an additive mathematical review of the complete supplied NOTE.md, equations (1)--(23), and its complete LaTeX counterpart. Both original files remain byte-for-byte unchanged. The proof of the Hankel criterion and the infinite-tail separator is accepted after the explicit symbol repair below. This proof supplies explicit growth constants, the inverse local coordinate, the complete finite algebra map, and the original theta-boundary representative of that map. The source note's numerical results retain their exploratory status.

The source defines \(u_N=\Re w_N-(\Re\beta)v_N\), but its boxed equation (16) uses the undefined symbol \(\nu_N\). The exact repair is the identity \[
\nu_N:=u_N=\Re w_N-(\Re\beta)v_N.
\] This occurs at NOTE.md lines 212/216 and NOTE.tex lines 350/354. The user's pasted narrative already uses \(u_N\) in both places. Equation (WBR16) below writes that complete definition and the polynomial using the same \(u_N\). No value or coefficient is changed.

\subsection{WBR1. Source and scalar factors}\label{C47stage:WBR:wbr1.-source-and-scalar-factors}

Let \[
D=-x\partial_x,\quad
\widehat\phi(\xi)=\int_{\mathbb R}\phi(x)e^{-2\pi ix\xi}\,dx,
\] \[
V=\{\phi\in\mathcal S(\mathbb R):\phi(-x)=\phi(x),\
 \phi(0)=0,\ \int_{\mathbb R}\phi=0\},
\] \[
\mathscr B=\{F\in C^\infty(\mathbb R_{>0}):
 \sup_{x>0}x^b|D^jF(x)|<\infty\ (b\in\mathbb Z,\ j\ge0)\},
\] \[
\Theta:V\to\mathscr B,\quad \Theta\phi(x)=2\sum_{n\ge1}\phi(nx),
\qquad \mathcal MF(s)=\int_0^\infty F(x)x^s\,\frac{dx}{x}.
\tag{WBR1}
\] The quotient is the original \(Q=\mathscr B/\Theta V\), with its quotient topology. Its split coefficient set is \(G(R)=\{\tau\}\sqcup R^\bullet\), with represented zero \(e_R=0_R^\bullet\). A map on a coefficient fibre sends \((\ell,v)\) to \((\ell,fv)\) and fixes external absence \(\tau\). Thus \((\ell,0)\) keeps the source label.

Put \[
q_0=e^{-\pi x^2},\quad
\phi_*=(D^2-D)q_0=(4\pi^2x^4-6\pi x^2)e^{-\pi x^2},
\quad f_0=\Theta\phi_*.
\tag{WBR2}
\] Direct differentiation gives \(Dq_0=2\pi x^2q_0\) and \(D^2q_0=(4\pi^2x^4-4\pi x^2)q_0\). The function \(\phi_*\) is even and vanishes at zero. Integration by parts gives \(\int Df=\int f\) on Schwartz functions, so \(\int\phi_*=0\). Hence \(\phi_*\in V\). Fourier differentiation and integration by parts give \(\mathcal FD=(1-D)\mathcal F\). Since \(\mathcal Fq_0=q_0\), this proves \(\mathcal F\phi_*=\phi_*\). Poisson summation, with zero origin terms on both sides, gives \[
f_0(x)=x^{-1}f_0(1/x).
\tag{WBR3}
\] For \(x\ge1\), the series and its derivatives decay faster than every power. Reflection proves the same assertion at zero, so \(f_0\in\mathscr B\).

For \(\Re s>1\), absolute convergence allows termwise integration: \[
\begin{aligned}
\mathcal Mf_0(s)
&=2\zeta(s)\int_0^\infty
 (4\pi^2x^4-6\pi x^2)e^{-\pi x^2}x^{s-1}\,dx\\
&=2\zeta(s)\pi^{-s/2}
 \{2\Gamma(s/2+2)-3\Gamma(s/2+1)\}\\
&=s(s-1)\pi^{-s/2}\Gamma(s/2)\zeta(s)=2\xi(s).
\end{aligned}
\tag{WBR4}
\] The two-ended decay makes the Mellin transform entire. Thus this identity also supplies the entire continuation of the completed zeta expression. Write \(g=2\xi\), retaining its actual value at \(1/2\). The conventional factor in \(\xi\) is given by \href{https://dlmf.nist.gov/25.4.E4}{DLMF 25.4.4}.

\subsection{WBR2. Positive density and explicit growth}\label{C47stage:WBR:wbr2.-positive-density-and-explicit-growth}

Define \[
\psi(y)=e^{y/2}f_0(e^y),\qquad
C_\theta=8\pi^2\sum_{n=1}^\infty n^4e^{-\pi(n^2-1)}.
\tag{WBR5}
\] The constant is finite. Reflection gives \(\psi(-y)=\psi(y)\). For \(y\ge0\), each theta summand is strictly positive because \(2\pi n^2e^{2y}>3\). Dropping its negative term only for an upper bound, while retaining the complete summand in the definition, yields \[
0<\psi(y)\le
8\pi^2e^{9y/2}\sum_{n\ge1}n^4e^{-\pi n^2e^{2y}}
\le C_\theta e^{9y/2-\pi e^{2y}}.
\tag{WBR6}
\] Reflection gives the same estimate with \(|y|\) on the whole line. For any compact set of \(s\), this dominates every \(|y|^j\psi(y)e^{(\Re s-1/2)y}\) by an integrable function. Therefore \[
g(s)=\int_{\mathbb R}\psi(y)e^{(s-1/2)y}\,dy,\qquad
M_{2j}=\int_{\mathbb R}y^{2j}\psi(y)\,dy=g^{(2j)}(1/2).
\tag{WBR7}
\] In particular \(g(s)>0\) for every real \(s\), without numerical input.

The entire function \(G(z)=g(1/2+iz)\) is even. Its Taylor series gives \[
G(z)=\sum_{j\ge0}\frac{(-1)^jM_{2j}}{(2j)!}z^{2j}
=F(z^2),\quad F(w)=\sum_{j\ge0}a_jw^j,\quad
a_j=\frac{(-1)^jM_{2j}}{(2j)!}.
\tag{WBR8}
\] The map \(z\mapsto z^2\) sends the circle of radius \(\sqrt R\) onto the circle of radius \(R\), so the convergent Taylor series proves that \(F\) is entire and \(M_F(R)=M_G(\sqrt R)\). No square-root branch is chosen. Its unchanged mass is \(F(0)=a_0=g(1/2)>0\).

Set \(a(R)=9/4+\sqrt R/2\). Substituting \(t=e^{2y}\), with \(dy=dt/(2t)\), into the majorant gives \[
\begin{aligned}
M_F(R)&\le2C_\theta\int_0^\infty
 e^{(9/2+\sqrt R)y-\pi e^{2y}}\,dy\\
&=C_\theta\int_1^\infty t^{a(R)-1}e^{-\pi t}\,dt\\
&\le C_\theta\pi^{-a(R)}\Gamma(a(R))\\
&\le2C_\theta\pi^{-a(R)}
 \left(\frac{2(a(R)-1)}{e}\right)^{a(R)-1}.
\end{aligned}
\tag{WBR9}
\] For the last inequality, maximize \(t^{a-1}e^{-t/2}\) at \(2(a-1)\) and integrate the remaining \(e^{-t/2}\). Here \(a(R)>1\). Thus \(\log M_F(R)=O(\sqrt R\log R)\), and the order is at most \(1/2\).

\subsection{WBR3. Reciprocal zeros and coefficient map}\label{C47stage:WBR:wbr3.-reciprocal-zeros-and-coefficient-map}

Use the standard genus-zero case of Hadamard factorization in this precise form: an entire function of finite order strictly below one, nonzero at zero, is its value at zero times the product \(1-w/\alpha\) over its zeros with multiplicity; the reciprocal absolute zero values are summable, and no nonconstant exponential factor occurs. WBR2 verifies its hypotheses for the actual \(F\).

Let \(\alpha_j\) be its distinct zeros, \(m_j\) their orders, and \(\beta_j=1/\alpha_j\). Then \[
F(w)=a_0\prod_j(1-\beta_jw)^{m_j},\qquad
S_\beta=\sum_jm_j|\beta_j|<\infty.
\tag{WBR10}
\] The reciprocal zeros are nonzero and have only zero as a possible accumulation point. On a disk where all \(|\beta_jw|<1/2\), the logarithmic derivative is dominated by \(2\sum m_j|\beta_j|\). Geometric expansion and absolute convergence give \[
-\frac{F'(w)}{F(w)}
=\sum_j\frac{m_j\beta_j}{1-\beta_jw}
=\sum_{n\ge0}b_nw^n,\qquad
b_n=\sum_jm_j\beta_j^{n+1}.
\tag{WBR11}
\] Multiplying by \(F\) and equating coefficients, with \(a_0\) retained, gives \[
b_n=\frac{-(n+1)a_{n+1}-\sum_{j=1}^na_jb_{n-j}}{a_0}.
\tag{WBR12}
\] These \(b_n\) are real because all \(a_j\) are real.

At a zero \(\rho=1/2+\delta+i\gamma\), the exact coordinates are \[
z=-i(\rho-1/2)=\gamma-i\delta,\quad \alpha=z^2,\quad
\beta=-\frac1{(\rho-1/2)^2}
=\frac{\gamma^2-\delta^2+2i\gamma\delta}
       {(\gamma^2+\delta^2)^2}.
\tag{WBR13}
\] Since \(g(1/2)>0\), \(z\ne0\), so \(z\mapsto z^2\) has nonzero derivative and preserves the local zero order. The two reflected points \(\rho,1-\rho\) supply one zero \(\alpha\) of that order. Conjugation preserves the orders and sends \(\alpha,\beta\) to their conjugates. For real \(v\), \(F(-v^2)=g(1/2-v)>0\). Thus no nonpositive real zero occurs, and RH is equivalent to every \(\beta_j\) being positive real. An off-line zero has \(\delta\ne0\), while positivity of \(g\) on the real axis gives \(\gamma\ne0\); its reciprocal coordinate is therefore nonreal by (WBR13).

\subsection{WBR4. Finite criterion and full infinite-tail proof}\label{C47stage:WBR:wbr4.-finite-criterion-and-full-infinite-tail-proof}

For \(P(X)=\sum_{r=0}^dc_rX^r\), with real coefficients, put \[
H_d=(b_{r+s})_{0\le r,s\le d},\qquad
\mathcal L(P^2)=c^TH_dc
=\sum_jm_j\beta_jP(\beta_j)^2.
\tag{WBR14}
\] The reciprocal zeros together with zero form a compact set. Thus boundedness of \(P\) and (WBR10) prove absolute convergence. Conjugate terms make the sum real. Under RH its terms are nonnegative, proving \(H_d\succeq0\). The square retained in this formula is the algebraic \(P(\beta)^2\).

For the converse, choose a nonreal reciprocal zero \(\beta\) of order \(m\), and its conjugate. Set \(r=|\beta|/2\). Among the finitely many distinct reciprocal zeros of modulus at least \(r\), remove these two and call the remaining set \(\mathcal A\). Define \[
B(X)=\prod_{\gamma\in\mathcal A}(X-\gamma),\qquad
B_0=\prod_{\gamma\in\mathcal A}(r+|\gamma|).
\tag{WBR15}
\] The empty products are one. Conjugation permutes \(\mathcal A\), so \(B\) is real and \(B(\beta)\ne0\). Choose either root \(c\) of \(\beta c^2=-1\). For an integer \(N\ge0\), put \[
w_N=\frac{c}{\beta^NB(\beta)},\quad
v_N=\frac{\Im w_N}{\Im\beta},\quad
u_N=\Re w_N-(\Re\beta)v_N,\quad
P_N(X)=X^NB(X)(u_N+v_NX).
\tag{WBR16}
\] The coefficients are real, \(u_N+v_N\beta=w_N\), and hence \(P_N(\beta)=c\) and \(P_N(\overline\beta)=\overline c\). The selected contributions total \(-2m\). Every other point of modulus at least \(r\) contributes zero, regardless of its multiplicity.

The remaining infinite tail has the explicit constants \[
C=\frac{|c|B_0}{|B(\beta)|}
 \left(1+\frac{|\Re\beta|+r}{|\Im\beta|}\right),\qquad
T=\sum_{|\beta_j|<r}m_j|\beta_j|.
\tag{WBR17}
\] Indeed \(|v_N|\le |w_N|/|\Im\beta|\) and \(|u_N|\le |w_N|(1+|\Re\beta|/|\Im\beta|)\). For \(|X|\le r\), \(|B(X)|\le B_0\), giving \(|P_N(X)|\le C2^{-N}\). Therefore \[
\left|\sum_{|\beta_j|<r}m_j\beta_jP_N(\beta_j)^2\right|
\le C^2T4^{-N},\qquad
\mathcal L(P_N^2)\le-2m+C^2T4^{-N}.
\tag{WBR18}
\] The tail is real by conjugation, which justifies the ordered inequality. If \(T=0\), take \(N=0\). If \(T>0\), the explicit choice \[
N=\max\!\left(0,\ 1+
 \left\lfloor\log_4\left(\frac{C^2T}{2m}\right)\right\rfloor\right)
\tag{WBR19}
\] satisfies \(4^N>C^2T/(2m)\). This yields a finite negative direction of degree at most \(N+|\mathcal A|+1\), proving \[
\mathrm{RH}\quad\Longleftrightarrow\quad
H_d\succeq0\quad\hbox{for every }d\ge0.
\tag{WBR20}
\] Every actual zero and every multiplicity is included in this proof. The degree depends on the specified pair and tail; no uniform degree or actual off-line zero is thereby supplied.

For a strict negative vector \(c\), set \(M=\max_{i,j}|(H_d)_{ij}|\) and \(n=d+1\). If \(\|\widetilde c-c\|_\infty\le\varepsilon\), expansion gives \[
|\widetilde c^TH_d\widetilde c-c^TH_dc|
\le M(2\|c\|_1n\varepsilon+n^2\varepsilon^2).
\tag{WBR21}
\] If the original value is \(-\mu<0\), choose \(\varepsilon\) making the right side less than \(\mu\), then choose rational coefficients within that distance. This proves the rational-witness assertion with the original matrix fixed. The shifted family has quadratic form \(\sum m_j\beta_j^2P(\beta_j)^2\), and is also nonnegative under RH.

\subsection{WBR5. Actual packet and its complete representative}\label{C47stage:WBR:wbr5.-actual-packet-and-its-complete-representative}

Let \(Z\) be a nonempty finite set of actual nontrivial zeros of \(g\), stable under reflection and conjugation, with every full order retained. Write \[
h(s)=\prod_{\rho\in Z}(s-\rho)^{m_\rho},\quad d_h=\deg h,\quad
E_h=\mathbb C[s]/(h),\quad A_h=M_s,\quad
\nu_h=g/h,\quad \eta_h=j_h(1/\nu_h).
\tag{WBR22}
\] The map \(j_h\) keeps all Taylor coefficients through orders \(m_\rho-1\). The element \(\eta_h\) is a unit because every order in \(h\) is the full order in \(g\). Let \(R_Z(u)\) be the unique polynomial of degree below \(d_h\) representing \(\eta_hu\).

The original Euler inverse on a zero Mellin value is \[
S_\rho F(x)=x^{-\rho}\int_x^\infty F(y)y^{\rho-1}\,dy,\qquad
(D-\rho)S_\rho F=F,\quad
\mathcal MS_\rho F(s)=\frac{\mathcal MF(s)}{s-\rho}.
\tag{WBR23}
\] The zero Mellin value makes the upper integral the negative lower integral. At infinity the bound \(F(x)\le Cx^{-N}\) yields \(C/(N-\Re\rho)\); at zero \(F(x)\le C'x^N\) yields \(C'/(N+\Re\rho)\), when the denominators are positive. These inequalities, and \(DS_\rho F=\rho S_\rho F+F\), prove every required seminorm bound. The homogeneous function \(x^{-\rho}\) is excluded by the two-ended decay, so the inverse is unique and \(D-\rho\) is injective. Iteration gives \(S_h\) on the vanishing full Mellin jets, independent of inversion order by uniqueness.

The actual section and observation are \[
R_{\rm ref}u=S_h\Theta(R_Z(u)(D)\phi_*),\quad
\sigma_hu=[R_{\rm ref}u]\in Q,\quad J_hF=j_h\mathcal MF,
\] \[
\mathcal MR_{\rm ref}u=\nu_hR_Z(u),\quad
J_hR_{\rm ref}=1_{E_h},\quad J_h\Theta=0.
\tag{WBR24}
\] The transform \(gR_Z(u)\) has all the zero jets needed by \(S_h\). The section identity is \(j_h\nu_h\,\eta_h=1\). Thus \(\sigma_h\) is injective. Polynomial division gives \[
DR_{\rm ref}-R_{\rm ref}A_h=f_0\ell_h,\qquad
\ell_h(u)=[s^{d_h-1}]R_Z(u).
\tag{WBR25}
\] Indeed \(sR_Z(u)-R_Z(su)=h(s)\ell_h(u)\); the two sides of (WBR25) have equal Mellin transforms. Mellin injectivity follows from Fourier injectivity in logarithmic coordinates. Passing to the original quotient proves \(D_Q\sigma_h=\sigma_hA_h\).

The source sequence cited by the package is \[
0\longrightarrow Q[h(D)]\longrightarrow E_h
\xrightarrow{\times j_hg}E_h
\longrightarrow Q/h(D)Q\longrightarrow0.
\tag{WBR26}
\] Its complete proof, including the source Euler inverse, full-jet surjectivity and the actual arrows, is included verbatim in dependencies/arithmetic\_input.tex, (A5)--(A10). Equations (WBR23)--(WBR25) also prove directly every packet inclusion and operator identity used by this witness, with the full \(\eta_h\).

\subsection{WBR6. Reciprocal action, inverse jet, and literal theta boundary}\label{C47stage:WBR:wbr6.-reciprocal-action-inverse-jet-and-literal-theta-boundary}

Since \(g(1/2)>0\), \(A_h-\tfrac12I\) is invertible. Define on this finite packet \[
T_h=-(A_h-\tfrac12I)^{-2}.
\tag{WBR27}
\] Through \(\sigma_h\), this is also an endomorphism of the actual finite image in \(Q\). The construction has this precise domain; no inverse on unrelated quotient classes is required.

At a primary block write \(\lambda_\rho=\rho-1/2\) and \(A_h=\rho I+N_\rho\), with full nilpotent length \(m_\rho\). The exact finite expansion is \[
T_h=\sum_{j=0}^{m_\rho-1}
 (-1)^{j+1}(j+1)\lambda_\rho^{-j-2}N_\rho^j.
\tag{WBR28}
\] This is the product \(-\lambda_\rho^{-2}(1+N_\rho/\lambda_\rho)^{-2}\). Multiplying the series by \((\lambda_\rho+N_\rho)^2\) gives constant \(-1\) and zero at every positive degree below \(m_\rho\), which proves the identity in the truncated algebra.

Put \(\beta_\rho=-\lambda_\rho^{-2}\) and \(Y_\rho=T_h-\beta_\rho I\). The coefficient of \(N_\rho\) is \(2\lambda_\rho^{-3}\ne0\). The complete inverse coordinate is \[
N_\rho=\lambda_\rho\sum_{j=1}^{m_\rho-1}
 \binom{-1/2}{j}(Y_\rho/\beta_\rho)^j.
\tag{WBR29}
\] This is a finite algebra identity without any analytic branch choice. In the truncated algebra the element \(U\) with constant term one and \(U^2(1+Y_\rho/\beta_\rho)=1\) is unique: each successive coefficient has leading scalar \(2\). The displayed binomial coefficients solve these recursive equations. The element \(1+N_\rho/\lambda_\rho\) solves the same equation since \(T_h/\beta_\rho=(1+N_\rho/\lambda_\rho)^{-2}\), proving (WBR29). For \(m_\rho=1\), the positive-degree sums are empty and \(N_\rho=0\). Thus every original local length is retained.

There is a literal original-source representative of this action. Let \(t_h(s)\) be the unique polynomial of degree below \(d_h\) representing \(-(s-1/2)^{-2}\) in \(E_h\). For a polynomial \(P\), let \(p_h(s)\) be the degree-below-\(d_h\) representative of \(P(t_h(s))\). Define \[
K_{P,h}(u;s)=
\frac{p_h(s)R_Z(u)(s)-R_Z(P(T_h)u)(s)}{h(s)}.
\tag{WBR30}
\] The numerator vanishes modulo \(h\) because multiplication by \(\eta_h\) commutes with \(P(T_h)\). Thus the quotient is a polynomial in the original variable \(s\). Mellin transformation gives the exact original boundary \[
p_h(D)R_{\rm ref}u-R_{\rm ref}P(T_h)u
=\Theta\bigl(K_{P,h}(u;D)\phi_*\bigr).
\tag{WBR31}
\] Both transforms are \(gK_{P,h}(u;s)\), proving the formula and every sign. It specifies the difference of representatives before passing to \(Q\), and proves the finite packet intertwining there. The support lift keeps the label of that theta boundary fixed.

\subsection{WBR7. Full finite algebra, half-trace, and scalar radical}\label{C47stage:WBR:wbr7.-full-finite-algebra-half-trace-and-scalar-radical}

Let \(\mathcal B_Z\) be the set of distinct reciprocal coordinates in this packet. A fibre of \(\rho\mapsto\beta_\rho\) consists exactly of the distinct pair \(\rho,1-\rho\): equality of the reciprocal squares gives equality or negation of \(\rho-1/2\). Their orders agree by reflection. Call this common order \(m_\beta\), and put \[
f_Z(X)=\prod_{\beta\in\mathcal B_Z}(X-\beta)^{m_\beta},
\qquad \mathcal A_Z=\mathbb C[X]/(f_Z).
\tag{WBR32}
\] The map \(P\mapsto P(T_h)\) has kernel exactly \((f_Z)\). On each primary block, the inverse (WBR29) identifies the local algebra with \(\mathbb C[Y]/Y^{m_\rho}\). Consequently the action vanishes precisely when every Taylor coefficient of \(P\) at \(\beta_\rho\) below that order is zero. Intersecting the coprime primary ideals proves the stated kernel and the embedding \(\mathcal A_Z=\mathbb C[T_h]\subset\operatorname{End}(E_h)\).

Choose one point of each reflected pair to write coordinates. Applying (WBR29) on both primary blocks and then Chinese remainders gives an explicit \(\mathcal A_Z\)-module isomorphism \[
E_h\simeq\mathcal A_Z\oplus\mathcal A_Z.
\tag{WBR33}
\] The action of \(X\) on each copy is multiplication by \(X\). Both reflected blocks remain in this display. Changing the coordinate choice swaps the corresponding two components without altering \(E_h\) or \(T_h\).

On the local basis \(1,Y,\ldots,Y^{m_\beta-1}\), multiplication by \(Q(\beta+Y)\) is triangular with constant diagonal \(Q(\beta)\). Therefore \[
\frac12\operatorname{Tr}_{E_h}Q(T_h)
=\operatorname{Tr}_{\mathcal A_Z}M_Q
=\sum_{\beta\in\mathcal B_Z}m_\beta Q(\beta),
\] \[
\frac12\operatorname{Tr}_{E_h}\bigl(T_hP(T_h)^2\bigr)
=\sum_{\beta\in\mathcal B_Z}m_\beta\beta P(\beta)^2.
\tag{WBR34}
\] This proves the original factor \(1/2\), the full multiplicity and the source-to-scalar map. Exhaustion by finite stable packets converges to (WBR14) by (WBR10).

The precise scalar information loss is also computable. On \(\mathcal A_Z\), define the bilinear form \[
\mathfrak b_Z(U,V)=\operatorname{Tr}_{\mathcal A_Z}M_{XUV}
=\sum_\beta m_\beta\beta U(\beta)V(\beta).
\tag{WBR35}
\] Its radical is exactly \[
\mathfrak n_Z=\{U:U(\beta)=0\text{ for every }\beta\}
             =(\operatorname{rad}f_Z)/(f_Z).
\tag{WBR36}
\] One inclusion follows from (WBR35). If \(U(\beta_0)\ne0\), choose a polynomial \(V\) equal to one at \(\beta_0\) and zero at every other node, by Lagrange interpolation. Then the bilinear value is \(m_{\beta_0}\beta_0U(\beta_0)\ne0\), proving the other inclusion. Thus the scalar form factors through the exact quotient \(\mathcal A_Z\to\prod_\beta\mathbb C\), whose kernel is the displayed nilradical. Before that quotient, (WBR28)--(WBR33) retain every nilpotent coefficient; the trace weights retain every algebraic multiplicity. This specifies the full map relating the local algebra to its scalar observation.

\subsection{WBR8. Exact critical-line comparison}\label{C47stage:WBR:wbr8.-exact-critical-line-comparison}

The inherited source map is \[
\Gamma:Q\longrightarrow Q_{\rm crit}
=\mathcal S(\mathbb R)/
\overline{\zeta(1/2+it)\mathcal S(\mathbb R)},\qquad
\Gamma[F]=[\mathcal MF(1/2+it)].
\tag{WBR37}
\] Its construction and range are included whole in dependencies/HOCHSCHILD\_COMPARISON.md, (H16)--(H25). For \(\Re\rho\ne1/2\), multiplication by \((1/2+it-\rho)^{-1}\) is a continuous Schwartz multiplier, preserves the displayed relation and its closure, and inverts \(M_{1/2+it}-\rho\) on \(Q_{\rm crit}\). Intertwining by \(D\) therefore sends every full generalized \(\rho\)-block to zero under \(\Gamma\sigma_h\).

At a critical zero \(\rho=1/2+i\gamma\) of full order \(m\), the continuous jet map \[
[f]\longmapsto
\sum_{j=0}^{m-1}\frac{i^{-j}f^{(j)}(\gamma)}{j!}\epsilon^j
\tag{WBR38}
\] vanishes on the relation and its closure. On \(\Gamma[R_{\rm ref}u]\), it equals \(j_\rho u\): each \(t\)-derivative gives the factor \(i^j\) in the Mellin derivative. Together these jet maps invert every critical block. Thus \[
\ker(\Gamma\sigma_h)=E_{h,\rm off}.
\tag{WBR39}
\]

In the reciprocal algebra, the corresponding actual map is the Chinese-remainder projection \[
\pi_{\rm crit}:\mathcal A_Z\longrightarrow
\mathcal A_{Z,\rm crit},\qquad
\ker\pi_{\rm crit}=\mathcal A_{Z,\rm off}.
\tag{WBR40}
\] The functional (WBR34) decomposes into its critical and off-critical summands. If an off-critical node \(\beta\) occurs, its primary idempotent \(e_\beta\), equal to one on its entire local block and zero on the others, belongs to \(\ker\pi_{\rm crit}\), yet has \(\operatorname{Tr}_{\mathcal A_Z}M_{Xe_\beta}=m_\beta\beta\ne0\). Hence in that case the full functional does not factor through \(\pi_{\rm crit}\). Its exact omitted part is the displayed weighted trace on \(\mathcal A_{Z,\rm off}\). On real polynomials the two-value interpolation (WBR16) selects the conjugate pair and kills every other finite packet node, giving the value \(-2m\); the actual infinite arithmetic tail is controlled by (WBR18).

All these maps lift with the source label fixed. Thus an off-critical vector maps to \((\ell,0)\) under the split lift of \(\Gamma\sigma_h\), while remaining the input \((\ell,v)\) of the full source operator and its trace. The exact maps (WBR29)--(WBR40) relate the arithmetic source, the scalar witness and the receiving quotient.

\subsection{WBR9. Computation and integration record}\label{C47stage:WBR:wbr9.-computation-and-integration-record}

The original integral in the numerical script is \[
M_{2j}=4\sum_{n\ge1}\int_0^\infty y^{2j}
(4\pi^2n^4e^{9y/2}-6\pi n^2e^{5y/2})
e^{-\pi n^2e^{2y}}\,dy.
\tag{WBR41}
\] The factor four is the two in \(\Theta\) and the two half-lines under reflection. The supplied script retains these constants. It uses finite theta sums and the interval \(0\le y\le4\). Its two high-precision runs have no validated enclosure of the finite quadrature or omitted tails.

The finite independent replay verifies the saved decimal matrices by exact rational LDL arithmetic. Those matrices are positive through dimension 16 in both shifts. This is a theorem about the specified decimal rational matrices. The map from those matrices to actual \(H_d\) requires coefficient enclosures; none is supplied by agreement between two floating runs.

The supplied verifier also permits negative shift values, using Python negative indices, and the coefficient recurrence uses floating division when passed integer inputs. These narrow input-contract defects do not affect the archived default runs or the mathematical recurrence. The additive strict verifier accepts explicit exact rationals, validates nonnegative integral indices and interval order, and leaves the archive unchanged. The full replay report, exact tests and machine receipts accompany this proof.

The source NOTE equations (1)--(23), with the explicit \(\nu_N=u_N\) symbol repair above, are accepted with their stated standard Hadamard input. For cumulative integration, include its entire original LaTeX, this entire proof, the complete finite replay and both inherited dependency files. Place the reciprocal action after the original arithmetic packet, and its Hankel scalar after the original theta moments. Its source class remains available to the ongoing metric and absolute-base geometric calculations.

Classical context was checked directly at \href{https://arxiv.org/abs/1510.03420}{Zhang's author record}. This lies in the established power-sum approach. The current proof makes no novelty claim for that general approach and does not depend on an uninspected theorem from that paper.
